\chapter{Recording, Verification, and Closure}

Every reading records the question, date, deck or instrument, method, shuffle
state, jumper policy, exact result, raw observations, interpretation, expansions,
uncertainty, and later events. A cross-system expansion also records its parent
object and its interpretive priority.

The method permits intuition, but it does not require memory to be perfect.
Writing the draw before interpreting it protects the reading from retrospective
rewriting. Later events are appended only after they occur.

The reader may close a reading when the question is answered, a practical action
is clear, the instrument signals rest, repetition begins, or the system reaches
its $T\_STOP$. More cards do not automatically produce more truth.

The aim is not to prove fate. It is to preserve a living instrument well enough
that its patterns, warnings, discoveries, and failures can be learned from.
